\documentclass[11pt]{article}
\usepackage{graphicx}
\usepackage{amssymb}
\usepackage{bm}
\usepackage{mathtools}
\usepackage{amsmath}
\usepackage{commath}
\usepackage{amsthm}
\usepackage{enumitem}
\usepackage{float} 
\usepackage{array}
\usepackage{outlines}
\usepackage{setspace}
\usepackage[colorinlistoftodos]{todonotes}
\definecolor{green}{rgb}{0.4, 0.69, 0.2}
\definecolor{asparagus}{rgb}{0.53, 0.66, 0.42}
\definecolor{blue}{rgb}{0.0, 0.53, 0.74}
\definecolor{terracotta}{HTML}{e2725b}
\definecolor{mocassin}{HTML}{FFE4B5}
\definecolor{coral}{HTML}{FFAA80}
\setlist[enumerate,1]{label=\roman*)}
\setlist[enumerate,2]{label=\alph*)}
\newtheorem{theorem}{Theorem}
\newtheorem{corollary}{Corollary}
\newtheorem{lemma}{Lemma}
\newtheorem{remark}{Remark}
\newtheorem{definition}{Definition}
\newtheorem{prop}{Proposition}
\newcommand*\diff{\mathop{}\!\mathrm{d}}
\usepackage{tikz}
\usetikzlibrary{arrows,shapes,trees, calc}
\usepackage{graphicx}
\usepackage{subfig}


%%% *** PREAMBLE *** %%%
\title{Diffusive velocities}

\date{}

%%%***  CONTENT *** %%%
\begin{document}

In this document \textit{advective velocity} means wind and essentially the advective air velocity, while \textit{c-diffusive velocity} stands for the velocities generated by the diffusion process of the contaminant ($c$). We use the term \textit{c-diffusive} to highlight that this is not the diffusion of air in itself, but the diffusion effect of the contaminant.

\begin{enumerate}

  \item In the section of strong solutions we exploit the fact that the velocity can be calculated without the concentration, or in other words, that while the $c$ does depend on $\mathbf{u},$ the velocity $\mathbf{u}$ does not depend on $c.$ For example, the $(\mathbf{u}, c)$ strong solution in the hydrostatic case is constructed in a certain order: we obtain it in a way that firstly we take $\mathbf{u}$ as a solution of the momentum equations, and then we find a $c$ accordingly. This means that $c$ can not have any effect or feedback on $\mathbf{u},$ otherwise it wouldn't make sense to find $\mathbf{u}$ without considering $c.$
  
  \item According to the previous point, the presence of the concentration function does not effect $\mathbf{u},$ thus the velocity $\mathbf{u}$ is exclusively advective in the sense that it is not affected by the diffusion process that takes place thanks to the presence of the contaminant: the c-diffusive part of air velocity is strictly zero.
  
  \item At the same time, we have $\mathcal{O}(1)$ diffusivity coefficients, which means that there is a significant diffusive flux of the contaminant, even without the presence of an outer, advective wind effect. This means that there is a non-zero diffusive velocity of the contaminants.
  
  \item Since we have non-zero c-diffusive contaminant velocities and zero c-diffusive air velocities, it follows that wind particles and contaminant particles do not have the same velocities. In other words, the $\mathbf{u}$ function we use throughout our work is the velocity of the air particles,
  
  $\mathbf{u} = \mathbf{u}^{\text{air}} = \mathbf{u}^{\text{air}}_{\text{advective}} + \mathbf{u}^{\text{air}}_{\text{c-diffusive}} = \mathbf{u}^{\text{air}}_{\text{advective}} + 0,$
  
  where $\mathbf{u}^{\text{air}}_{\text{c-diffusive}} = 0$ according to point \textit{ii)}.
  
  At the same time there exists the velocity of the contaminants (we don't use a notation for it, nor is it used in the thesis), which is non-zero thanks to diffusion effects even in the case of no wind advection.
  
\end{enumerate}

\textbf{Conclusion}:

We do not use the assumption of air velocity and contaminant velocity being the same.

\end{document}

